\documentclass[10pt]{scrartcl}
\usepackage[a4paper, total={6in, 8in}]{geometry}
\usepackage[usenames,dvipsnames,svgnames]{xcolor}
\usepackage[shortlabels]{enumitem}
\usepackage[framemethod=TikZ]{mdframed}
\usepackage{amsmath,amssymb,amsthm}
\usepackage[colorlinks]{hyperref}

\hypersetup{
	colorlinks=true,
	linkcolor=blue,
	filecolor=blue,      
	urlcolor=blue,
	pdftitle={MAT 487 Spr25}
}

\usepackage{microtype}
\usepackage{mathtools}
\usepackage[headsepline]{scrlayer-scrpage}
\usepackage{thmtools}
\usepackage[textsize=small]{todonotes}

\addtolength{\textheight}{3.14cm}
\providecommand{\ol}{\overline}
\providecommand{\eps}{\varepsilon}
\providecommand{\half}{\frac{1}{2}}
\providecommand{\dang}{\measuredangle} %% Directed angle
\providecommand{\CC}{\mathbb C}
\providecommand{\FF}{\mathbb F}
\providecommand{\NN}{\mathbb N}
\providecommand{\QQ}{\mathbb Q}
\providecommand{\RR}{\mathbb R}
\providecommand{\ZZ}{\mathbb Z}
\providecommand{\dg}{^\circ}
\providecommand{\ind}[1]{\mathbf 1_{#1}}
\providecommand{\ii}{\item}
\providecommand{\alert}[1]{{\sffamily\textbf{\textcolor{magenta}{#1}}}}
\providecommand{\opname}{\operatorname}
\providecommand{\li}{\opname{li}}
\providecommand{\ls}{\opname{ls}}
\providecommand{\inv}{^{-1}}
\providecommand{\setdiff}{\smallsetminus}
\providecommand{\mb}[1]{\mathbf{#1}}
\providecommand{\abs}[1]{\left|{#1}\right|}
\providecommand{\symdiff}{\Delta}
\providecommand{\jim}{m_{*,(J)}}
\providecommand{\jom}{m^{*,(J)}}
\providecommand{\pcint}{\mathrm{p.c.}\int}
\providecommand{\dif}{\;d}
\providecommand{\dist}{\operatorname{dist}}
\providecommand{\diam}{\operatorname{diam}}

\reversemarginpar
\setlength{\marginparwidth}{3cm}

\mdfdefinestyle{mdbluebox}{roundcorner=10pt,innerbottommargin=9pt,
	linecolor=blue,backgroundcolor=TealBlue!5,}
\declaretheoremstyle[headfont=\sffamily\bfseries\color{MidnightBlue},
mdframed={style=mdbluebox},]{thmbluebox}

\mdfdefinestyle{mdredbox}{frametitlefont=\bfseries,innerbottommargin=8pt,
	nobreak=true,backgroundcolor=Salmon!5,linecolor=RawSienna,}
\declaretheoremstyle[headfont=\bfseries\color{RawSienna},
mdframed={style=mdredbox},headpunct={\\[3pt]},postheadspace=0pt,]{thmredbox}

\mdfdefinestyle{mdgreenbox}{linecolor=ForestGreen,backgroundcolor=ForestGreen!5,
	linewidth=2pt,rightline=false,leftline=true,topline=false,bottomline=false,}
\declaretheoremstyle[headfont=\bfseries\sffamily\color{ForestGreen!70!black},
mdframed={style=mdgreenbox},headpunct={ --- },]{thmgreenbox}

\mdfdefinestyle{mdorangebox}{linecolor=RedOrange,backgroundcolor=RedOrange!5,
	linewidth=2pt,rightline=false,leftline=true,topline=false,bottomline=false,}
\declaretheoremstyle[headfont=\bfseries\sffamily\color{RedOrange!70!black},
mdframed={style=mdorangebox},headpunct={ --- },]{thmorangebox}

\mdfdefinestyle{mdblackbox}{linecolor=black,backgroundcolor=RedViolet!5!gray!5,
	linewidth=3pt,rightline=false,leftline=true,topline=false,bottomline=false,}
\declaretheoremstyle[mdframed={style=mdblackbox}]{thmblackbox}


\declaretheorem[style=thmredbox,name=Problem,numberwithin=section]{problem}

\declaretheorem[style=thmredbox,name=Example,numbered=no]{example}
\declaretheorem[style=thmredbox,name=$\bigstar$ Problem,sibling=problem]{niceprob}

\declaretheorem[style=thmbluebox,name=Theorem,sibling=problem]{theorem}
\declaretheorem[style=thmbluebox,name=Corollary,sibling=theorem]{corollary}
\declaretheorem[style=thmbluebox,name=Proposition,sibling=theorem]{prop}
\declaretheorem[style=thmbluebox,name=Lemma,sibling=theorem]{lemma}
\declaretheorem[style=thmbluebox,name=Theorem,numbered=no]{theorem*}
\declaretheorem[style=thmbluebox,name=Lemma,numbered=no]{lemma*}
\declaretheorem[style=thmgreenbox,name=Claim,numberwithin=theorem]{claim}
\declaretheorem[style=thmgreenbox,name=Claim,numbered=no]{claim*}
\declaretheorem[style=thmblackbox,name=Remark,sibling=theorem]{remark}
\declaretheorem[style=thmblackbox,name=Remark,numbered=no]{remark*}
\declaretheorem[style=thmgreenbox,name=Definition,sibling=theorem]{definition}
\declaretheorem[style=thmgreenbox,name=Definition,numbered=no]{definition*}
\declaretheorem[style=thmorangebox,name=Warning,sibling=theorem]{warning}
\declaretheorem[style=thmorangebox,name=Warning,numbered=no]{warning*}
\declaretheorem[style=thmorangebox,name=Note,numbered=no]{note}
\declaretheorem[style=thmblackbox,name=Exercise,numberwithin=section]{exercise}
\declaretheorem[style=thmblackbox,name=Exercise,numbered=no]{exercise*}

\newenvironment{soln}{\begin{proof}[Solution]}{\end{proof}}
\newenvironment{walkthrough}{\noindent \textbf{Walkthrough.}}{\medskip}
\newlist{walk}{enumerate}{3}
\setlist[walk]{label=\bfseries (\alph*)}

\providecommand{\pow}{\operatorname{Pow}}
\providecommand{\vocab}[1]{{\textbf{\textcolor{ForestGreen}{#1}}}}
\providecommand{\lcm}{\operatorname{lcm}}
\providecommand{\ord}{\operatorname{ord}}
\providecommand{\ex}{\opname{ex}}
\providecommand{\floor}[1]{\left\lfloor #1\right\rfloor}
\providecommand{\ceil}[1]{\left\lceil #1\right\rceil}

\usepackage{asymptote}
\begin{asydef}
	defaultpen(fontsize(10pt));
	unitsize(1cm);
	size(8cm); // set a reasonable default
	usepackage("amsmath");
	usepackage("amssymb");
	settings.tex="pdflatex";
	settings.outformat="pdf";
	// Replacement for olympiad+cse5 which is not standard
	import geometry;
	// recalibrate fill and filldraw for conics
	void filldraw(picture pic = currentpicture, conic g, pen fillpen=defaultpen, pen drawpen=defaultpen)
	{ filldraw(pic, (path) g, fillpen, drawpen); }
	void fill(picture pic = currentpicture, conic g, pen p=defaultpen)
	{ filldraw(pic, (path) g, p); }
	// some geometry
	pair foot(pair P, pair A, pair B) { return foot(triangle(A,B,P).VC); }
	pair orthocenter(pair A, pair B, pair C) { return orthocentercenter(A,B,C); }
	pair centroid(pair A, pair B, pair C) { return (A+B+C)/3; }
	// cse5 abbreviations
	path CP(pair P, pair A) { return circle(P, abs(A-P)); }
	path CR(pair P, real r) { return circle(P, r); }
	pair IP(path p, path q) { return intersectionpoints(p,q)[0]; }
	pair OP(path p, path q) { return intersectionpoints(p,q)[1]; }
	path Line(pair A, pair B, real a=0.6, real b=a) { return (a*(A-B)+A)--(b*(B-A)+B); }
	// cse5 more useful functions
	picture CC() {
		picture p=rotate(0)*currentpicture;
		currentpicture.erase();
		return p;
	}
	pair MP(Label s, pair A, pair B = plain.S, pen p = defaultpen) {
		Label L = s;
		L.s = "$"+s.s+"$";
		label(L, A, B, p);
		return A;
	}
	pair Drawing(Label s = "", pair A, pair B = plain.S, pen p = defaultpen) {
		dot(MP(s, A, B, p), p);
		return A;
	}
	path Drawing(path g, pen p = defaultpen, arrowbar ar = None) {
		draw(g, p, ar);
		return g;
	}
\end{asydef}

\ihead{\footnotesize\textbf{MAT 487 Spring 2025}}
\ohead{\footnotesize \today}

\title{\vspace{-2.0cm} Tur\'an-flavored stuff}
\author{\large Aditya Pahuja}
\date{\large Presented on 2025-02-05}
\begin{document}
\maketitle
\tableofcontents
\pagebreak

All graphs here are simple.

\section{Mantel's theorem and Tur\'an's theorem}
Our goal throughout this lecture will be to think about the following question:
\begin{problem}
	For a positive integer $n$ and graph $H$,
	let \vocab{$\ex(n, H)$} be the largest possible number of edges
	in an $n$-vertex graph that does not contain $H$ as a subgraph
	(such graphs are called \vocab{$H$-free}).
	How does $\ex(n, H)$ grow as $n\to\infty$?
\end{problem}

To start with, let's consider a very simple choice of $H$:
\begin{problem}
	What is $\ex(n, K_3)$, the largest possible number of edges
	in an $n$-vertex triangle-free graph?
\end{problem}

In the early 1900s, Willem Mantel proved Mantel's theorem, namely that
\[\ex(n, K_3) = \floor{\frac{n^2}{4}}.\]
We shall present two proofs that $\floor{n^2/4}$ is an upper bound
on the number of edges.

\begin{proof}[Proof 1]
	Let $G$ be a triangle-free graph with $m$ edges.
	Let $xy$ be an edge of $G$ with endpoints $x$ and $y$.
	Then, the set $N(x)$ (the vertices connected to $x$)
	is disjoint from $N(y)$, meaning
	\[\deg x + \deg y \le n\]
	for all $x,y\in V$. Summing over all edges in $G$ gives
	\[\sum_{xy\in E} \deg x + \deg y \le mn.\]
	Since each vertex $x$ appears in $\deg x$ edges, the summation equals
	\[\sum_{x\in V}(\deg x)^2.\]
	Applying the Cauchy-Schwarz inequality then gives
	\[mn \ge \sum_{x\in V}(\deg x)^2 \ge
	\frac 1n\left(\sum_{x\in V} \deg x\right)^2 = \frac{(2m)^2}{n},\]
	which simplifies to $m \le n^2/4$, but we can take the floor on the right
	because $m$ is an integer.
\end{proof}

\begin{proof}[Proof 2]
	Let $v$ be a vertex of maximal degree, let $A = N(v)$, and let $B$
	be the vertices that are not in $A$.
	Then, there are no edges connecting two vertices of $A$
	(else we would have a triangle),
	meaning every single edge in $G$ has an endpoint in $B$. Thus
	\[m \le \sum_{x\in B} \deg x
	\le |B|\cdot \max_{x\in B}\deg x = |B|\deg v = |A|\cdot|B|,\]
	but since we know that $|A| + |B| = n$ is fixed, we can apply
	the AM-GM inequality to get $m\le |A|\cdot|B|\le n^2/4$,
	at which point taking floors finishes.
\end{proof}

To show that this bound on $\ex(n, K_3)$ is sharp, we need to find
a graph where equality holds. To do this, let's trace through the proofs
to see what properties a maximal graph would need to preserve
equality the whole way.

In the first proof, we would need $\deg x + \deg y = n$
for all edges for the $mn$ upper bound to be sharp.
This happens iff $G$ is a complete bipartite graph on $n$ vertices.
From here, it's easy to check that the complete bipartite graph
that maximizes the number of edges is $K_{\floor{n/2}, \ceil{n/2}}$.

Similarly, in the second proof,
\[m = \sum_{x\in B}\deg x\]
can only hold if $G$ is bipartite, and then the upper bound on this sum
requires the graph to be complete bipartite, from which we
recover the same unique maximal graph.

\begin{exercise}
	Go do some of the problems in section 1.1 of the book.
\end{exercise}

One next step is to look at arbitrary cliques $K_{r+1}$.
\begin{problem}
	What is $\ex(n, K_{r+1})$?
\end{problem}

Analogous to Mantel's theorem, we can immediately see that
the $n$-vertex $r$-partite graph with part sizes differing by at most $1$,
which we will call the \vocab{Tur\'an graph} $T_{n,r}$,
is $K_{r+1}$-free. It turns out that this is again the best we can do:

\begin{theorem}[Tur\'an's theorem]
	The Tur\'an graph is the unique maximal $K_{r+1}$-free graph, and
	\[\ex(n, K_{r+1})\le \left(1 - \frac1r\right)\frac{n^2}{2}.\]
	Moreover, this bound is asymptotically tight
	for fixed $r$ as $n\to\infty$:
	\[\ex(n, K_{r+1}) = \left(1 - \frac1r - o(1)\right)\frac {n^2}2.\]
\end{theorem}

\begin{proof}[Proof 1]
	This proof mirrors the first proof of Mantel's theorem.
	We will induct on $r$, with $r = 0$ being trivial.
	
	To start with, let $v$ again be a vertex of maximal degree
	in a $K_{r+1}$-free graph $G$, and again let $A = N(v)$ and
	$B = V\setdiff A$. Then, $A$ is $K_r$-free (else the $K_r$ along with
	$B$ makes a $K_{r+1}$), so $e(A)\le e(T_{|A|, r})$.
	On the other hand, the number of edges that have an endpoint in $B$
	can again be bounded by
	\[\sum_{x\in B}\deg x\le |B|\cdot\max_{x\in B}\deg x = |A|\cdot|B|,\]
	implying that equality happens when $G$ is complete $r$-partite.
	
	To show that the Tur\'an graph is the specific $r$-partite graph
	that maximizes the number of edges, we can just observe that
	if two parts of the graph have counts of vertices that differ by
	more than $1$, then moving a vertex from the more populous part
	into the less populous one will increase the number of edges.
\end{proof}




















	
\end{document}